\section{連続な関数は積分できる}\label{節:積分可能性}
定義は積分が存在する場合の意味を定めたが，存在そのものをまだ保証していない。本書で主に扱う連続関数については，次の定理を用いる。

\begin{theorem}{連続関数の積分可能性：採用する事実}\label{定理:連続積分可能}
閉区間 $[a,b]$ で連続な関数は，その区間でリーマン積分可能である。
\end{theorem}

この定理の完全な証明は，本書では扱わない。考え方は，十分短い小区間では関数の高さの変動も小さくなり，どの代表点を選んでも和の差が小さくなるというものである。閉区間上の連続性を使うと，その小ささを区間全体でそろえられる。各小区間での最大の高さと最小の高さの差が $\eta$ 以下なら，上から作る和と下から作る和の差は
\[
 \sum_{k=1}^{n}\eta\Delta x_k=\eta(b-a)
\]
以下になる。このような評価をすべての分割について保証する部分も含め，上の定理として採用する。

したがって，連続関数の定積分の値を計算するときは，存在をこの定理で確かめたうえで，計算しやすい等分割と代表点を選んでよい。右端を選ぶなら，
\[
 \int_a^b f(x)\dd x
 =\lim_{n\to\infty}\frac{b-a}{n}
  \sum_{k=1}^{n}f\left(a+\frac{k(b-a)}n\right)
\]
となる。ここでは幅が $\frac{b-a}{n}$，第 $k$ 小区間の右端が $a+\frac{k(b-a)}n$ である。

\begin{example}{リーマン和から値を求める}
区間 $[1,3]$ を $n$ 等分し，各小区間の右端を使って $\int_1^3(2x-1)\dd x$ を求めよ。
\end{example}
被積分関数は連続なので，定積分は存在する。幅は $\frac2n$，右端は $1+\frac{2k}{n}$ である。したがって，リーマン和は
\begin{align*}
 S_n
 &=\frac2n\sum_{k=1}^{n}\left\{2\left(1+\frac{2k}n\right)-1\right\}\\
 &=\frac2n\sum_{k=1}^{n}\left(1+\frac{4k}n\right)\\
 &=2+\frac8{n^2}\cdot\frac{n(n+1)}2
  =6+\frac4n.
\end{align*}
よって $n\to\infty$ として，定積分は $6$ となる。この例では和の公式が使えたが，一般の関数では和を直接計算するのは難しい。後で微分との関係を使い，別の計算法を導く。

